\chapter{曲线系}
\section{曲线系引入}
\begin{example}[（椭圆上的中点弦直线）]
设椭圆$\dfrac{x^2}{a^2}+\dfrac{y^2}{b^2}=1$内有一点$P(x_0,y_0)$，求以$P$为弦中点的直线方程.
\end{example}
\begin{solution}
    设弦上一端为 $X(x,y)$。由于 $P(x_0,y_0)$ 是中点，另一端为 $X'(2x_0-x,2y_0-y)$。两端都在椭圆上，故联立
    \[\begin{cases}\dfrac{x^2}{a^2}+\dfrac{y^2}{b^2}=1\\\dfrac{(2x_0-x)^2}{a^2}+\dfrac{(2y_0-y)^2}{b^2}=1\end{cases}\]
    两式相减，得\[\dfrac{x_0x}{a^2}+\dfrac{y_0y}{b^2}=\dfrac{x_0^2}{a^2}+\dfrac{y_0^2}{b^2}\]
\end{solution}
\begin{example}[（椭圆上的直线两点式）]
设椭圆$\dfrac{x^2}{a^2}+\dfrac{y^2}{b^2}=1$上有两点$A(x_1,y_1),B(x_2,y_2)$，求$AB$方程.
\end{example}
\begin{solution}
    构造同样经过 $A,B$ 的二次式，再与原方程相减，可得到过 $A,B$ 的一次方程。取
    \[\dfrac{(x-x_1)(x-x_2)}{a^2}+\dfrac{(y-y_1)(y-y_2)}{b^2}=0\]
    与原椭圆方程相减，整理得$AB$的方程
    \[\dfrac{x_1+x_2}{a^2}x+\dfrac{y_1+y_2}{b^2}y=1+\dfrac{x_1x_2}{a^2}+\dfrac{y_1y_2}{b^2}\]
    代入椭圆参数方程，可写成参数形式的两点式。
\end{solution}
\begin{example}[（抛物线上的直线两点式）]
设$y^2=2px$上有两点$A(x_1,y_1),B(x_2,y_2)$，求$AB$方程.
\end{example}
\begin{solution}
    取同过 $A,B$ 的直线对 $(y-y_1)(y-y_2)=0$，与抛物线方程相减，得
    \[y^2-(y-y_1)(y-y_2)=2px\Leftrightarrow 2px-(y_1+y_2)y+y_1y_2=0\]
\end{solution}
\begin{example}[（利用梯形得到四点共圆）]
    二次函数$y=x^2+(a+1)x+a-2$与坐标轴交于$A,B,C$三点，求$\triangle ABC$的外接圆的方程，并验证圆心是否在定直线上，以及外接圆是否过定点.
\end{example}
\begin{solution}
    抛物线与坐标轴的交点中，两点在 $x$ 轴上，另一点在 $y=a-2$ 这条水平线上。于是用 $y(y-a+2)=0$ 表示经过这两条水平线的退化二次式，再与抛物线
    \[
    x^2+(a+1)x-y+a-2=0
    \]
    相加，得
    \[x^2+y^2+(a+1)x-y-ay+2y+a-2=x^2+y^2+(a+1)x+(1-a)y+a-2=0\]
    圆心为
    \[
    \left(-\frac{a+1}{2},\frac{a-1}{2}\right),
    \]
    其横纵坐标之和恒为 $-1$，所以圆心一定在定直线
    \[
    x+y+1=0
    \]
    上。将方程改写为
    \[a(x-y+1)+x^2+y^2+x+y-2=0\]
    代入$(0,1)$和$(-2,-1)$时，含 $a$ 的项都为 $0$，所以外接圆恒过定点$(0,1),(-2,-1)$。
\end{solution}
    \begin{example}[（2025八省联考）]
    椭圆$\dfrac{x^2}{9}+\dfrac{y^2}{8}=1$，$A(-3,0),B(2,0)$，过$B$的弦$MN$与点$A$形成的三角形$AMN$外心为$D$，证明$k_{MN}k_{OD}$为定值.
\end{example}
\begin{solution}
    若弦 $MN$ 竖直，则 $k_{MN}$ 不存在，与题意不符；故设弦 $MN$ 的方程为 $x-ty-2=0$。要得到 $\triangle AMN$ 的外接圆，需要先构造一个过 $A,M,N$ 的二次式。取过 $A(-3,0)$ 的直线 $x+ty+3=0$，与 $MN$ 合成
\[\begin{cases}(x-ty-2)(x+ty+3)=0\\8x^2+9y^2-72=0\end{cases}\]
    前一式已经过 $A,M,N$。再取原椭圆方程的适当倍数，使 $x^2,y^2$ 的系数相等，就得到过 $A,M,N$ 的圆：
    \begin{align*}&(t^{2}+1)(8x^{2}+9y^{2}-72)+(x-ty-2)(x+ty+3)=0\\&\Leftrightarrow{}(8t^{2}+9)x^{2}+(8t^{2}+9)y^{2}+x-5ty-72t^{2}-78=0\end{align*}
    则$k_{OD}=-5t,k_{MN}=\dfrac1{t}$，定值为$-5$.
\end{solution}
\begin{example}[（直径圆的求法）]
    已知抛物线$y^2=2px$和弦$AB$所在的直线$x=ty+m$，求$A,B$的直径圆方程.
\end{example}
\begin{solution}
    消去$x$得到\[y^2=2p(ty+m)=2pty+2pm\Rightarrow y^2-2pty-2pm=0\]
    消去$y$得到\[(x-m)^2=t^2y^2=2pt^2x\Rightarrow x^2-(2m+2pt^2)x+m^2=0\]
    两式相加，得直径圆方程
    \[x^2+y^2-(2m+2pt^2)x-2pty+m^2-2pm=0\]
    上述相加的依据是：消去$x,y$后分别得到关于两端点横坐标、纵坐标的二次方程，两式相加即得到以$AB$为直径的圆。
\end{solution}
\begin{example}[（过两点的圆系方程例题）]
    已知$C:y^2=4x,l:x=1$交于$A,B$两点，求经过$A,B$且与$x=-1$相切的圆的方程.
\end{example}
\begin{solution}
    联立 $x=1$ 与 $y^2=4x$，得
    \[
    A(1,2),\qquad B(1,-2).
    \]
    过 $A,B$ 的圆心必在 $x$ 轴上，设圆方程为
    \[
    x^2+y^2+Dx+F=0.
    \]
    代入 $A(1,2)$，得
    \[
    5+D+F=0.
    \]
    该圆与直线 $x=-1$ 相切，等价于把 $x=-1$ 代入后关于 $y$ 的方程有重根：
    \[
    y^2+1-D+F=0.
    \]
    因为其中没有一次项，所以重根只能为 $y=0$，故
    \[
    1-D+F=0.
    \]
    解得 $D=-2,\ F=-3$。因此所求圆为
    \[
    x^2+y^2-2x-3=0.
    \]
\end{solution}
\begin{example}[（直径圆习题）]
    已知椭圆$3x^2+4y^2=12$，弦$AB$过定点$(m,0)$，其直径圆与椭圆交于点$C,D$，证明直线$CD$过定点.
\end{example}
\begin{solution}
    设$AB:x=ty+m$，其中$m$已知。分别消去 $x$ 与 $y$，由韦达定理得到端点横、纵坐标满足的二次方程：
    \[\begin{cases}3x^{2}+4y^{2}-12=0\\x=ty+m&\end{cases}\Rightarrow\begin{cases}(3t^{2}+4)y^{2}+6tmy+3m^{2}-12=0\\(3t^{2}+4)x^{2}-8mx+4m^{2}-12t^{2}=0\end{cases}\]
    两式相加得到:\[(3t^2+4)x^2+(3t^2+4)y^2-8mx+6tmy+7m^2-12-12t^2=0\]
    减去适当倍数的原椭圆方程，可把它整理成两个一次因子的乘积：
    \[-(t^2+1)(3x^{2}+4y^{2}-12)+(3t^{2}+4)(x^{2}+y^{2})-8mx+6tmy+7m^{2}-12-12t^{2}=0\]
    因式分解成\[(x-ty-m)(x+ty-7m)=0\]，所以$CD$过定点$(7m,0)$.
\end{solution}
\begin{example}[（角平分线）]
    已知椭圆$3x^2+4y^2=12$，弦$AB$过定点$P(1,0)$，$C(4,0)$，证明$x$轴平分$\angle{ACB}$
\end{example}
\begin{solution}
    由$AB:x=ty+1$。设 $A'B'$ 为 $AB$ 关于 $x$ 轴的对称直线，则 $A'B':x=-ty+1$，对应端点也是 $A,B$ 关于 $x$ 轴的对称点。于是得到直线对
    \[(x-ty-1)(x+ty-1)=0\]
    再减去适当倍数的原椭圆方程。为了让所得式子表示两条过 $C(4,0)$ 的对称直线，把常数项配成 $(x-4)^2$ 的形式，由此确定倍数为 $-\dfrac14$：
    \[-(3x^{2}+4y^{2}-12)+4(x-ty-1)(x+ty-1)=0\]
    得到$(x-4)^2=(4+t^2)y^2$，这等价于\[(x+\sqrt{4+t^2}y-4)(x-\sqrt{4+t^2}y-4)=0\]
    这两条直线分别经过 $A,B$，斜率互为相反数，所以$x$轴平分$\angle ACB$。
\end{solution}
\newpage
\begin{example}[（焦点、外心与面积）]
    椭圆$x^2+4y^2=4$的右焦点为$(\sqrt3,0)$，点$M,N$在$x$轴上方的椭圆上，且$\triangle{MNF}$的外心$D$在$x$轴上，$\sqrt{3}|FM||FN|=2|MN|$，求$S_{\triangle{MNF}}$
\end{example}
\begin{solution}
    设直线方程$MN:y=kx+m\Leftrightarrow kx=y-m$，并联立
    \vspace{-5pt}
    \[\begin{cases}(1):(4k^2+1)x^2+8kmx+4m^2-4=0\\(2):(4k^2+1)y^2-2my+m^2-4k^2=0\\(3):y-kx-m=0\end{cases}\vspace{-5pt}\]
    将$(1)+(2)+2m(3)$相加，得过$M,N$且圆心在$x$轴上的圆\vspace{-5pt}
    \[(4k^2+1)(x^2+y^2)+6kmx+3m^2-4k^2-4=0\vspace{-5pt}\]
    代入$F(\sqrt3,0)$，得直线$MN$到焦点$(\sqrt3,0)$的距离 $d$：\vspace{-5pt}
    \[8k^2+6\sqrt{3}km+3m^2=1\Leftrightarrow d=\dfrac{\sqrt{3}k+m}{\sqrt{k^2+1}}=\dfrac{1}{\sqrt3}\vspace{-5pt}\]
    设圆半径为$R$。由等面积关系，
    \vspace{-5pt}\[\begin{cases}S_{\triangle{MNF}}=\dfrac12|FM||FN|\sin\angle{MFN}=\dfrac12|MN|d=\dfrac{1}{2}\dfrac{1}{\sqrt3}|MN|\\\sqrt{3}|FM||FN|=2|MN|\end{cases}\Rightarrow \sin\angle{MFN}=\dfrac12\vspace{-5pt}\]
    设$M(x_1,y_1),N(x_2,y_2)$。由 $\sin\angle MFN=\frac12$，可设
    \[
    \cos\angle MFN=\varepsilon\frac{\sqrt3}{2},\qquad \varepsilon=\pm1.
    \]
    椭圆 $\frac{x^2}{4}+y^2=1$ 的右焦点为 $F(\sqrt3,0)$，右准线为 $x=\frac4{\sqrt3}$，故椭圆上点 $X(x,y)$ 到右焦点的距离为
    \[
    |FX|=2-\frac{\sqrt3}{2}x.
    \]
    于是
    \begin{align*}
        \dfrac{\overrightarrow{FM}\cdot\overrightarrow{FN}}{|FM||FN|}
        =&\dfrac{(x_1-\sqrt3)(x_2-\sqrt3)+y_1y_2}{(2-\frac{\sqrt3}{2}x_1)(2-\frac{\sqrt3}{2}x_2)}=\dfrac{(x_1-\sqrt3)(x_2-\sqrt3)+y_1y_2}{\frac34(\frac{4}{\sqrt3}-x_1)(\frac{4}{\sqrt3}-x_2)}\\
        =&\dfrac{4}{3}\dfrac{12k^2+8\sqrt{3}km+4m^2-1+m^2-4k^2}{(4k^2+1)\frac{16}{3}+\frac{32}{\sqrt3}km+4m^2-4}=\dfrac{8k^2+8\sqrt{3}km+5m^2-1}{16k^2+8\sqrt{3}km+3m^2+1}\\
        =&\frac{8k^2+8\sqrt{3}km+5m^2-(8k^2+6\sqrt{3}km+3m^2)}{16k^2+8\sqrt{3}km+3m^2+8k^2+6\sqrt{3}km+3m^2}=\frac{\sqrt3km+m^2}{12k^2+7\sqrt3km+3m^2}
    \end{align*}
    若 $\varepsilon=-1$，令 $t=\frac{m}{k}$ 后得到
    \[
    (2+3\sqrt3)t^2+(21+2\sqrt3)t+12\sqrt3=0,
    \]
    其两根为 $-\sqrt3$ 与 $-\dfrac{12}{2+3\sqrt3}$。代入
    \[
    k^2=\frac{d^2}{(t+\sqrt3)^2-d^2}
    \]
    时，分母均为负数，故 $k^2<0$，不合直线斜率为实数。故只能取 $\varepsilon=1$。令$t=\dfrac{m}{k}$，把式子同除以$k^2$，得
    \[(3\sqrt3-2)t^2+(21-2\sqrt3)t+12\sqrt3=0.\]
    两根中，$t=-\sqrt3$ 会使 $\sqrt3k+m=0$，与前面 $d=\frac1{\sqrt3}$ 不符，因此取另一根
    \[t=\dfrac{12}{2-3\sqrt3},\qquad t^2=\dfrac{144}{31-12\sqrt3}.\vspace{-5pt}\]
    由距离式
    \[
    \frac{\sqrt3k+m}{\sqrt{k^2+1}}=\frac1{\sqrt3}
    \]
    两边平方，并把 $m=tk$ 代入，得到
    \[
    k^2=\frac{d^2}{(t+\sqrt3)^2-d^2},\qquad d=\frac1{\sqrt3}.
    \]
    于是
    \[k^2=\dfrac{d^2}{(-\frac{m}{k}-\sqrt3)^2-d^2}=\dfrac{1}{3(\frac{3+2\sqrt3}{3\sqrt3-2})^2-1}=\dfrac{31-12\sqrt3}{3(3+2\sqrt3)^2-(3\sqrt3-2)^2}=\dfrac{1}{16}\dfrac{31-12\sqrt3}{2+3\sqrt3}\vspace{-5pt}\]
    \[m^2=k^2\dfrac{m^2}{k^2}=\dfrac{1}{16}\dfrac{31-12\sqrt3}{2+3\sqrt3}\dfrac{144}{31-12\sqrt3}=\dfrac{9}{2+3\sqrt3},km=\dfrac34\dfrac{2-3\sqrt3}{2+3\sqrt3}\]\vspace{-10pt}
    于是
    \begin{align*}
        S_{\triangle{MNF}}&=\dfrac12\sin\angle{MFN}|FM||FN|=\dfrac{1}{4}|FM||FN|=\dfrac{16k^2+8\sqrt{3}km+3m^2+1}{4k^2+1}\\
        &=\dfrac{(31-12\sqrt3)+6\sqrt{3}(2-3\sqrt3)+27+2+3\sqrt3}{\frac14(31-12\sqrt3)+2+3\sqrt3}=\dfrac{2+\sqrt3}{13}
    \end{align*}
\end{solution}
\newpage
\begin{example}[（焦点、外心与面积，圆系解法）]
    椭圆$x^2+4y^2=4$的右焦点为$(\sqrt3,0)$，点$M,N$在$x$轴上方的椭圆上，且$\triangle{MNF}$的外心$D$在$x$轴上，$\sqrt{3}|FM||FN|=2|MN|$，求$S_{\triangle{MNF}}$
\end{example}
\begin{solution}
    也可以联立圆与椭圆。设外心为 $D(m,0)$，圆过焦点 $F(\sqrt3,0)$，故
    \[
    r^2=(m-\sqrt3)^2.
    \]
    于是
    \[\begin{cases}(x-m)^2+y^2=r^2\\\dfrac{x^2}{4}+y^2=1\end{cases}\Rightarrow \dfrac{3}{4}x^2-2mx+2\sqrt{3}m-2=0\]
    设$M(x_1,y_1),N(x_2,y_2)$，由$\sqrt{3}|FM||FN|=2|MN|$得到：
    \begin{align*}|MN|=&\dfrac{\sqrt3}{2}\left(2-\dfrac{\sqrt3}{2}x_1\right)\left(2-\dfrac{\sqrt3}{2}x_2\right)=\dfrac{3\sqrt3}{8}\left(\dfrac{4}{\sqrt3}-x_1\right)\left(\dfrac{4}{\sqrt3}-x_2\right)\\=&\dfrac{3\sqrt3}{8}\dfrac43\left(\dfrac{3}{4}\dfrac{16}{3}-\dfrac{8m}{\sqrt3}+2m-2\right)
    =\dfrac{\sqrt3}{2}(2-\dfrac{2}{\sqrt{3}}m)=\sqrt3-m\end{align*}
    由于圆经过椭圆焦点，有\[(\sqrt3-m)^2=r^2\Rightarrow \sqrt3-m=|MN|=r=\dfrac{\sqrt{3}}{2}|FM||FN|\]
    设圆心为$D$，则 $DM=DN=r$，又已得 $MN=r$，所以 $\triangle MDN$ 为等边三角形。点 $F$ 也在该圆上，圆周角 $\angle MFN$ 与圆心角 $\angle MDN$ 对同一弦 $MN$，因此 $\angle{MFN}=\dfrac{\pi}{6}$，所求面积为
    \[S_{\triangle{MNF}}=\dfrac12\sin\angle{MFN}|FM||FN|=\dfrac{1}{4}|FM||FN|=\dfrac{\sqrt3}{6}r\]
    $r$ 可由三角形$MNF$的余弦定理求出：
    \[|MF|^2+|NF|^2-\sqrt3|MF||NF|=|MF|^2+|NF|^2-2r=|MN|^2=r^2\Rightarrow |MF|^2+|NF|^2=r^2+2r\]
    再利用恒等式$|MF|^2+|NF|^2+2|MF||NF|=(|MF|+|NF|)^2$，结合韦达定理得：
    \[|MF|+|NF|=4-\dfrac{\sqrt3}{2}(x_1+x_2)=4-\dfrac{\sqrt3}{2}\dfrac{8}{3}m=4-\dfrac{4}{\sqrt3}m\]
    那么\[|MF|^2+|NF|^2=(|MF|+|NF|)^2-2|MF||NF|=(4-\dfrac{4}{\sqrt3}m)^2-\dfrac{4}{\sqrt3}r=r^2+2r\]
    代入$\sqrt3-m=r$得到\[\left(4-\dfrac{4}{\sqrt3}(\sqrt3-r)\right)^2-\dfrac{4}{\sqrt3}r=r^2+2r\Rightarrow r=\dfrac{6+4\sqrt3}{13}\]
    得到$S_{\triangle{MNF}}=\dfrac{\sqrt3}{6}r=\dfrac{2+\sqrt3}{13}$
\end{solution}

